\section{PRZYGOTOWANIE DO GRY}

Ogólna zasada:

Każdy z graczy powinien zaopatrzyć się w ołówek oraz jedną kartę cech dla każdej postaci, którą gra.

Żetony powinny zostać starannie porozdzielane oraz posegregowane według rodzajów.
Zaleca się, aby wszystkie zapisy na kartach cech dokonywane były ołówkiem, tak, aby łatwo je było można zetrzeć, gdyż w trakcie gry cechy postaci ulegają ciągłym zmianom.

Zasady szczegółowe:

[4.1] Przygotowanie do gry musi przebiegać w następujący sposób:

A. WYBÓR BOHATERÓW.

1. Postacie są podzielone na dwie grupy: bohaterowie oraz uczniowie.
Na żetonach, reprezentujących bohaterów znajduje się imię, zaś żetony, reprezentujące uczniów określają tylko ich rasę.

Żetony, reprezentujące bohaterów powinny zostać odwrócone zadrukowaną stroną ku dołowi oraz wymieszane tak, aby wybór był całkowicie przypadkowy.

2. Jeżeli gra tylko jeden gracz wybiera on trzech bohaterów (przypadkowo!) oraz trzech uczniów (dowolnie).

3. Jeżeli w grze uczestniczy więcej niż jedna osoba kolejność w jakiej dokonywany jest wybór bohaterów ustalana jest przy pomocy rzutu kostką.
Wybór bohaterów dokonywany jest kolejno, zgodnie z ruchem wskazówek zegara, przy czym jako pierwszy dokonuje wyboru (,,ciągnie'') ten gracz, który wyrzucił największą liczbę oczek.

4. Gracze ,,ciągną'' kolejno, aż do rozdzielenia między siebie trzech żetonów bohaterów i trzech żetonów uczniów.
Liczba postaci (tzn. bohaterów i uczniów) w grze może być mniejsza niż 6 (sugerowane minimum wynosi 4), jednak nie powinna przekraczać tej liczby.

B. OKREŚLENIE CECH POSTACI.

1. Gracze odszukują cechy swoich bohaterów w tabeli 4.3 ,,Charakterystyka bohaterów'' i zapisują uzyskane w tej tabeli informacje na karcie cech danego bohatera.

2. Gracze wybierają rasę swoich uczniów.
Uczeń może być krasnoludem, człowiekiem lub elfem.

3. W zależności od rasy na karcie cech ucznia zapisywane są następujące informacje:

TABELA

4. Gracze wybierają dwa rodzaje broni dla każdego ze swoich uczniów (dowolnie z tabeli 9.9).

5. Gracze rzucają 1K6 i z tabeli 4.4 ,,Potencjał magiczny'' odczytują wielkość współczynnika potencjał magiczny dla każdego ucznia, a następnie zapisują go na karcie cech.
Każdy współczynnik składa się z trzech cyfr.

6. Dla każdej postaci, której potencjał magiczny (jeden z jego trzech członów) jest większy od zera gracze wybierają czary, którymi ta postać będzie dysponować w czasie gry.
Liczba czarów musi być równa największemu z trzech członów, składających się na potencjał magiczny.
Na przykład: postać, posiadająca potencjał magiczny równy 4/3/2 powinna rozpocząć grę znając 4 różne czary.

7. Zakłada się, że każdy uczeń zdobył już w swoim dotychczasowym życiu pewne doświadczenie.
W związku z tym gracz może:
a. zwiększyć jego siłę o 1 lub
b. zwiększyć biegłość w jego podstawowej broni o 1 lub
c. zwiększyć jego zdolności o 1.

C. USTALENIE SZYKU MARSZOWEGO.

1. Raczej rzadko będzie się zdarzać, że postacie będą poruszać się po labiryncie samotnie.
Zakłada się, że idą one razem, tworząc grupę, która w instrukcji nazywana będzie oddziałem.
Dlatego dla oznaczenia miejsca, w którym postacie się znajdują używany jest tylko jeden żeton, noszący nazwę ,,Oddział''.

Jednak w każdej chwili powinny być znane pozycje, jakie postacie zajmują względem siebie, w ramach oddziału.
Służy temu ustalenie przed rozpoczęciem gry szyku marszowego.

2. Szyk marszowy składa się z pewnej liczby szeregów, przy czym w żadnym z nich nie może być więcej niż trzy postacie.
Gracze układają żetony, oznaczające postacie w dogodnym miejscu na stole tak, aby tworzyły one szyk marszowy.
Od decyzji graczy zależy która postać będzie w jakim szeregu i na jakiej pozycji.
W pierwszym szeregu szyku muszą być co najmniej dwie postacie (o ile aż tyle pozostaje przy życiu).

3. Gracze mogą dowolnie zmieniać szyk marszowy do chwili rozpoczęcia walki.
W jej trakcie szyk może zostać zmieniony tylko podczas fazy ,,reorganizacja oddziału'' (patrz 9.8).

D. Ustalenie dominującego słońca.

Jeden z graczy rzuca 1K6 w celu ustalenia które z trzech słońc jest dominujące czyli pod wpływem którego słońca oddział się znajduje.

Wpływ słońca określa liczbę czarów, które mogą być używane przez daną postać podczas bieżącej przygody (czyli w ciągu części gry liczonej od momentu wejścia do labiryntu do momentu wyjścia z niego w celu wyleczenia ran, zmiany broni, itp.).
Dla każdej kolejnej przygody należy na nowo ustalać dominujące słońce.
Rezultaty rzutu kostką odczytuje się następująco:

1 lub 2 oczka – dominuje czerwone słońce, liczba czarów określona jest lewą cyfrą współczynnika potencjał magiczny,

3 lub 4 oczka – dominuje żółte słońce – środkowa cyfra,

5 lub 6 oczek – dominuje niebieskie słońce – prawa cyfra.

Po ustaleniu dominującego słońca gracze wybierają z tych czarów, które zostały przypisane danej postaci w punkcie B6 taką ich liczbę, aby odpowiadała wielkości potencjału magicznego dla dominującego słońca.
Na przykład: jeżeli wyrzucone zostały 3 oczka, a potencjał magiczny postaci równy jest 4/5/6 gracz powinien wybrać (i zaznaczyć, np. przez podkreślenie) 5 spośród 6 posiadanych przez tę postać czarów.
Szósty, odrzucony czar podczas tej przygody nie będzie mógł być użyty.

E. SEGREGACJA ŻETONÓW.

Wszystkie żetony, reprezentujące komnaty i korytarze powinny zostać umieszczone w dużej filiżance lub innym tego typu pojemniku.
Żetony ,,Czarne Wrota'', ,,Brama labiryntu'' oraz ,,Bezimienny'' powinny zostać odłożone na bok.
Ich użycie jest określone specjalnymi przepisami.
Podobnie należy odłożyć na bok niewykorzystane żetony, reprezentujące bohaterów i uczniów.
W ciągu tej gry nie będą one już potrzebne.

F. ROZPOCZĘCIE GRY.

Żeton ,,Brama labiryntu'' kładziony jest na stole, oznaczając tym samym miejsce, w którym znajduje się wejście.
Na tym żetonie powinien zostać położony żeton ,,Oddział''.
W tym momencie gra może się rozpocząć.
Gracze przystępują do wykonywania czynności, podanych w części 5.0 ,,Przebieg gry''.

[4.2] Karta cech postaci – patrz plansza z tabelami.

[4.3] Tabela ,,Charakterystyka bohaterów'' – patrz plansza z tabelami.

[4.4] Tabela ,,Potencjał magiczny'' – patrz plansza z tabelami.
